%%=============================================================================
%% Voorwoord
%%=============================================================================

\chapter*{\IfLanguageName{dutch}{Woord vooraf}{Preface}}%
\label{ch:voorwoord}

%% TODO:
%% Het voorwoord is het enige deel van de bachelorproef waar je vanuit je
%% eigen standpunt (``ik-vorm'') mag schrijven. Je kan hier bv. motiveren
%% waarom jij het onderwerp wil bespreken.
%% Vergeet ook niet te bedanken wie je geholpen/gesteund/... heeft

Het schrijven van deze bachelorproef markeert het einde van mijn academische traject binnen de opleiding Toegepaste Informatica aan HOGENT. Het was een boeiende en uitdagende periode waarin ik mijn kennis en vaardigheden in informatica, data-analyse en business intelligence stap voor stap heb kunnen verdiepen en toepassen in praktijkgerichte projecten. Deze bachelorproef vormt het sluitstuk daarvan en focust op het ontwerpen en realiseren van een data-gedreven scoutingomgeving voor de BNXT League.

Basketbal speelt al jaren een belangrijke rol in mijn leven. Ik heb zelf twaalf jaar lang basketbal gespeeld en ben daarbij altijd gefascineerd geweest door statistieken en het verhaal achter de cijfers. Die combinatie van persoonlijke interesse en technische achtergrond maakte het voor mij een logische keuze om een onderwerp te kiezen waarin basketbal en data-analyse samenkomen. Met deze bachelorproef wil ik aantonen hoe een goed opgezet datawarehouse en een doordacht Power BI-dashboard scoutingbeslissingen kunnen ondersteunen en objectiveren.

Graag wil ik mijn promotor, de heer Stefaan Samyn, bedanken voor zijn waardevolle begeleiding en constructieve feedback gedurende het hele proces. Zijn kritische vragen en suggesties hebben mee vorm gegeven aan de opzet van het onderzoek en de uitwerking van het proof of concept.

Deze bachelorproef is niet alleen een afsluiting van mijn studie, maar ook een startpunt voor verdere ontwikkeling in het brede domein van sportdata en analytics. Ik hoop dat dit werk kan tonen dat er binnen competities zoals de BNXT League nog veel potentieel is voor datagedreven scouting.
